
\chapter{Modelli Finanziari in Spazi di Probabilità Finiti}

In questo capitolo esploriamo i fondamenti della matematica finanziaria in tempo discreto. L'obiettivo principale è mostrare come la teoria delle martingale costituisca lo strumento fondamentale per prezzare (valutare) i derivati finanziari e descrivere l'assenza di arbitraggio.

\section{Modelli uni-periodali (One-period models)}

Consideriamo un mercato finanziario elementare che vive solo in due istanti di tempo: $n=0$ (oggi) e $n=1$ (domani). Lo spazio di probabilità $(\Omega, \mathcal{F}, \PP)$ è finito: $\Omega = \{\omega_1, \dots, \omega_M\}$, con probabilità $\PP(\{\omega_k\}) = p_k > 0$.

Il mercato è composto da due asset:
\begin{itemize}
    \item Un asset privo di rischio (\textbf{Bond} o Conto in banca) $B$:
    \[ B_0 = 1, \quad B_1 = 1 + r \]
    dove $r > -1$ è il tasso di interesse (deterministico).
    \item Un asset rischioso (\textbf{Stock} o Azione) $S$:
    \[ S_0 = s > 0, \quad S_1(\omega) = s_k > 0 \text{ se } \omega = \omega_k \]
    Il valore $S_1$ è una variabile aleatoria (il prezzo di domani non è noto oggi).
\end{itemize}

È comodo introdurre i \emph{valori scontati} degli asset (prendendo il bond come \textit{numeraire}):
\[ B^*_n = \frac{B_n}{(1+r)^n} \implies B^*_0 = 1, \; B^*_1 = 1 \]
\[ S^*_n = \frac{S_n}{(1+r)^n} \implies S^*_0 = s, \; S^*_1 = \frac{S_1}{1+r} \]

\subsection{Portafogli e Arbitraggio}

Una \textbf{strategia di portafoglio} è un vettore $h = (h^B, h^S) \in \R^2$ che indica le quantità di bond e stock possedute. Il valore del portafoglio agli istanti $n=0, 1$ è:
\[ V_0^h = h^B B_0 + h^S S_0 = h^B + h^S s \]
\[ V_1^h = h^B B_1 + h^S S_1 = h^B (1+r) + h^S S_1 \]
Anche per il portafoglio definiamo il valore scontato: $V_n^{*h} = V_n^h / (1+r)^n = h^B B^*_n + h^S S^*_n$.

\begin{definition}[Arbitraggio]
Un \textbf{arbitraggio} è una strategia di portafoglio $h$ tale che:
\begin{enumerate}
    \item $V_0^h = 0$ (investimento iniziale nullo);
    \item $\PP(V_1^h \ge 0) = 1$ (nessun rischio di perdere denaro);
    \item $\PP(V_1^h > 0) > 0$ (probabilità strettamente positiva di fare profitto).
\end{enumerate}
Se non esistono strategie con queste proprietà, il mercato si dice \textbf{non arbitraggiabile} (arbitrage-free).
\end{definition}

\begin{hint}[Spiegazione Semplice dell'Arbitraggio]
Un arbitraggio è, letteralmente, una "macchina per fare soldi gratis". Entri nel mercato con zero euro (ad esempio andando corto sul bond per finanziare l'acquisto di stock), sei certo di non perdere mai nulla e hai una reale possibilità di uscirne più ricco. In un mercato efficiente, simili opportunità vengono immediatamente sfruttate (e quindi distrutte) dagli investitori. Pertanto, l'assenza di arbitraggio è il requisito minimo per un modello finanziario sensato.
\end{hint}

\subsection{Misure Martingala e Teoremi Fondamentali}

Per collegare i prezzi all'assenza di arbitraggio, introduciamo il concetto chiave del capitolo.

\begin{definition}[Misura Martingala Equivalente]
Una probabilità $\QQ$ su $\Omega$ si dice \textbf{equivalente} a $\PP$ (scritto $\QQ \sim \PP$) se hanno gli stessi eventi impossibili ($\QQ(A)=0 \iff \PP(A)=0$).
Una misura equivalente $\QQ$ è detta \textbf{Misura Martingala Equivalente (EMM)} se il prezzo scontato dello stock $S^*$ è una martingala sotto $\QQ$:
\[ \E^\QQ[S^*_1] = S^*_0 \iff \E^\QQ[S_1] = s(1+r) \]
\end{definition}

Sotto la probabilità $\QQ$, il rendimento atteso dell'asset rischioso è pari al rendimento dell'asset privo di rischio. $\QQ$ è nota anche come probabilità \emph{risk-neutral}.

\begin{theorem}[Primo Teorema Fondamentale dell'Asset Pricing]
Il mercato è non arbitraggiabile (arbitrage-free) \textbf{se e solo se} esiste (almeno) una misura martingala equivalente $\QQ$.
\end{theorem}

Un'altra questione fondamentale è la "completezza" del mercato. Un \textbf{Contingent Claim} (es. un'opzione) è un contratto che paga $\phi(S_1)$ al tempo 1. Il claim è detto \textbf{replicabile (o hedgeable)} se esiste un portafoglio $h$ tale che $V_1^h = \phi(S_1)$ per ogni possibile stato $\omega$. 
Se \emph{tutti} i possibili claim sono replicabili, il mercato si dice \textbf{completo}.

\begin{theorem}[Secondo Teorema Fondamentale dell'Asset Pricing]
Assumendo che il mercato sia non arbitraggiabile, esso è \textbf{completo} se e solo se esiste \textbf{esattamente una} misura martingala equivalente $\QQ$.
\end{theorem}

In questo modello a singolo stock, se $S_1$ può assumere 2 valori ($M=2$, modello binomiale) il mercato è completo. Se assume $M > 2$ valori, il mercato è incompleto (ci sono infinite misure $\QQ$).

\section{Pricing (Valutazione dei derivati)}

Se consideriamo un claim $\phi(S_1)$ in un mercato arbitrage-free, qual è il suo "prezzo giusto" (o \textit{fair price}) $x$ al tempo $n=0$?
Il prezzo $x$ deve essere tale che se inseriamo il derivato nel mercato come nuovo asset da scambiare, il mercato allargato rimanga privo di arbitraggi.

Abbiamo due strade (che coincidono per claim replicabili):
\begin{enumerate}
    \item \textbf{Hedging (Replica):} Se il claim è replicabile dal portafoglio $h$, allora per evitare arbitraggi il prezzo \emph{deve} essere il costo di replica oggi:
    \[ x = V_0^h \]
    Se fosse $x > V_0^h$, venderemmo il claim e compreremmo il portafoglio $h$, incassando la differenza subito a rischio zero. Viceversa se $x < V_0^h$.
    \item \textbf{Risk-Neutral Pricing:} In base al Primo Teorema, per non avere arbitraggi, anche il prezzo scontato del derivato deve essere una martingala sotto $\QQ$:
    \[ x = \E^\QQ\left[ \frac{\phi(S_1)}{1+r} \right] \]
\end{enumerate}

\begin{tcolorbox}[colback=red!5!white,colframe=red!75!black,title=Attenzione: Pricing sotto $\PP$ vs $\QQ$]
L'errore più comune nei modelli finanziari è calcolare il prezzo di un'opzione scontando il suo valore atteso rispetto alla probabilità oggettiva (o fisica) $\PP$. \textbf{Questo è sbagliato e porta ad arbitraggi!} Il prezzo di non-arbitraggio \emph{deve} essere calcolato usando il valore atteso sotto la misura risk-neutral $\QQ$. Sotto $\QQ$ gli investitori non richiedono alcun premio per il rischio.
\end{tcolorbox}

\section{Modelli multi-periodali: Il modello Binomiale (CRR)}

Estendiamo l'orizzonte a $N$ periodi di tempo: $n=0, 1, \dots, N$. Il mercato ha un bond deterministico $B_n = (1+r)^n$. L'asset rischioso $S_n$ evolve stocasticamente.

Una strategia $h = (h_n^B, h_n^S)_{n=0}^{N}$ è ora un \emph{processo pre-vedibile} (decido $h_n$ all'istante $n-1$, in base alle informazioni disponibili). La strategia è \textbf{autofinanziante} (self-financing) se le variazioni del valore di portafoglio sono dovute unicamente all'andamento del mercato, senza immissioni/prelievi di liquidità:
\[ h_{n+1}^B B_n + h_{n+1}^S S_n = h_n^B B_n + h_n^S S_n = V_n^h \]
Passando ai valori scontati:
\[ V_{n+1}^{*h} - V_{n}^{*h} = h_{n+1}^S (S_{n+1}^* - S_n^*) \]
che è esattamente una trasformata di martingala discreta! Ecco perché se $S^*$ è martingala (sotto $\QQ$), anche l'evoluzione di qualsiasi portafoglio autofinanziato $V^{*h}$ sarà una martingala sotto $\QQ$.

\subsection{Il Modello Binomiale di Cox-Ross-Rubinstein}
Nel modello binomiale, il prezzo dello stock passa dall'istante $n$ a $n+1$ moltiplicandosi per un fattore $Z_{n+1}$ che può assumere solo due valori: $u$ (up) e $d$ (down).
\[ S_{n+1} = Z_{n+1} S_n, \quad Z_{n+1} \in \{u, d\}, \quad 0 < d < u \]
Per evitare arbitraggi banali (in cui $S$ rende sempre più o sempre meno del bond), imponiamo:
\[ d < 1+r < u \]

Sotto questa ipotesi, l'equazione per la martingalità del prezzo scontato, ovvero $\E^\QQ[Z_{n+1}] = 1+r$, ci forza a determinare univocamente le probabilità risk-neutral di salita e discesa in ogni singolo passo temporale:
\[ q_u = \QQ(Z = u) = \frac{(1+r) - d}{u - d} \in (0, 1) \]
\[ q_d = \QQ(Z = d) = \frac{u - (1+r)}{u - d} \in (0, 1) \]
Essendo la misura $\QQ$ unica, il Primo e Secondo Teorema garantiscono che il Modello Binomiale sia \textbf{arbitrage-free} e \textbf{completo}.

\subsection{Prezzatura backward}
Data un'opzione europea con payoff $\phi(S_N)$ a scadenza $N$, il suo prezzo ad un generico istante $n$ è dato dal valore atteso condizionato:
\[ X_n = \frac{1}{(1+r)^{N-n}} \E^\QQ[\phi(S_N) \mid \mathcal{F}_n] \]
Nel modello binomiale, questa formula si esprime con la \emph{backward recursion}. Al tempo finale $X_N = \phi(S_N)$. Per un passo all'indietro:
\[ X_n(s) = \frac{1}{1+r} \left( q_u X_{n+1}(su) + q_d X_{n+1}(sd) \right) \]
Costruendo l'albero recombining dei prezzi possibili (binomial tree), possiamo partire dalle foglie (istante $N$) e lavorare a ritroso (istante per istante) calcolando per ogni nodo il valore pesato $X_n$ fino ad arrivare a $X_0$, che sarà il fair price al tempo iniziale.
